% !PRINTED 22
\weberpage{22}
From this, however, one easily recognises that these groups are not
only of the same degree, but that they are \gesperrt{isomorphic}. For we have
simultaneously
\[
Rb_1\,Rb_2 = Rb_1b_2
\]
and
\[
Q'b_1\,Q'b_2 = Q'b_1b_2.
\]
In the same way one can also conclude that the two groups
\[
Q'/R,\quad P/Q
\]
are isomorphic.

But the groups $P/Q$ and $P/Q'$ are, since $Q$, $Q'$ are greatest
normal divisors, simple by 3), and consequently the groups $Q'/R$
and $Q/R$ isomorphic with them are also simple, and consequently we
have the theorem:

\gesperrt{4. If $Q$ and $Q'$ are two different greatest normal
divisors of $P$, and if $R$ is the intersection of $Q$ and $Q'$,
then $R$ is a greatest normal divisor both of $Q$ and of $Q'$, and
the index of $R$ with respect to $Q$ is equal to the index of $Q'$
with respect to $P$.}

With this we have gained the foundation for demonstrating very
simply the theorem of the constancy of the index series.

For a better survey we write the various composition series to be
drawn into comparison in such a way that we place the index of each
term with respect to the preceding one under the term.

Accordingly, let any two given composition series of $P$, with the
associated indices, be the following:
\weq{23}{\begin{aligned}
&P,\ Q,\ Q_1,\ Q_2.\ Q_3 \ldots \\ % [sic: the source prints a full stop after Q_2.]
&\hphantom{P,\ }\nu,\ \nu_1,\ \nu_2,\ \nu_3 \ldots
\end{aligned}}
\weq{24}{\begin{aligned}
&P,\ Q',\ Q_1',\ Q_2',\ Q_3' \ldots \\
&\hphantom{P,\ }\nu',\ \nu_1',\ \nu_2',\ \nu_3' \ldots
\end{aligned}}
Now let $R$ be the intersection of $Q$ and $Q'$, and let $\mu$ and
$\mu'$ be the indices of $R$ with respect to $Q$ and $Q'$. On
account of the isomorphism of the groups $P/Q$ and $Q'/R$ we then
have $\mu' = \nu$, and for the corresponding reason $\mu = \nu'$. We
form a composition series of $R$ with the associated index series
\weq{25}{\begin{aligned}
&R,\ R_1,\ R_2,\ R_3 \ldots \\
&\hphantom{R,\ }\mu_1,\ \mu_2,\ \mu_3 \ldots,
\end{aligned}}
%CONT>>
